\subsection{More on enumerations}
%\phantom{+}%%\bigskip

 Since $\dif_{k+1}:\mcb_{k+1}\to\mcc_{k}$ is injective, it is
natural to consider in $\dif_{k+1}(\mcb_{k+1})$ the enumeration
inherited by the $\srle$ enumeration of $\mcb_{k+1}$ (see
Corollary~\ref{ltbij}, $(b)$). This is not a minor matter as regards
this paper. Indeed, to choose a ``correct'' enumeration in $\mcb_{k}$,
and hence in $G_{k-1}=\dif_{k}(\mcb_{k})$, has been seen to be crucial
in the proof of Lemma~\ref{superfluous} and hence in the proof of the
main theorem, Theorem~\ref{main}. In the light of this comment, the
next remark (and Remark~\ref{nonsrle} below) may help to clarify
ideas.

\begin{remark}\label{enum}
Let us consider in $\dif_{k+1}(\mcb_{k+1})$ the enumeration inherited
by $\mcb_{k+1}$. On the other hand, let us also consider in
$\Upsilon:=\bigcup_{i=2}^{r_k}\{\tau_{k}(e_{k,i},e_{k,j})\mid
e_{k,j}\in\mcb_{k,i}\}$ the enumeration defined in
Remark~\ref{induced-srle} (where $\Upsilon$ is possibly larger than
the set of syzygies considered in Remark~\ref{induced-srle}). Then
both enumerations agree, considering that
$-\dif_{k+1}(\mcb_{k+1})=\Upsilon$ (see the proof of
Theorem~\ref{main}).
\end{remark}
\begin{proof}
Let $e_{k+1,l}=(I_1,\ldots ,I_{k+2})\in\mcb_{k+1}$ and
$\dif_{k+1}(e_{k+1,l})= -\tau_{k}(e_{k,i},e_{k,j})$, with
\begin{eqnarray*}
e_{k,i}=(I_1,\ldots ,I_{k},I_{k+1}\cup I_{k+2})\mbox{ and }
e_{k,j}=(I_1,\ldots ,I_{k-1},I_k\cup I_{k+1},I_{k+2}).
\end{eqnarray*}
Analogously, let $e_{k+1,h}=(H_1,\ldots ,H_{k+2})\in\mcb_{k+1}$ and
$\dif_{k+1}(e_{k+1,h})=-\tau_{k}(e_{k,p},e_{k,q})$, with
\begin{eqnarray*}
e_{k,p}=(H_1,\ldots ,H_k,H_{k+1}\cup H_{k+2})\mbox{ and
}e_{k,q}=(H_1,\ldots H_{k-1},H_k\cup H_{k+1},H_{k+2}).
\end{eqnarray*}
(See Proposition~\ref{tau}.) Suppose that $\dif_{k+1}(e_{k+1,h})$
precedes $\dif_{k+1}(e_{k+1,l})$ in the inherited enumeration by
$\mcb_{k+1}$, that is, $e_{k+1,h}$ precedes $e_{k+1,l}$. By
Remark~\ref{srle}, there exists an integer $s\in\{1,\ldots ,k+1\}$,
such that $H_1=I_1,\ldots ,H_{s-1}=I_{s-1}$ and $H_s$ precedes
$I_s$. If $s\leq k$, then $e_{k,p}$ precedes $e_{k,i}$. If $s=k+1$,
then $e_{k,p}=e_{k,i}$ and $e_{k,q}$ precedes $e_{k,j}$. Therefore,
$\tau_k(e_{k,p},e_{k,q})$ precedes $\tau_k(e_{k,i},e_{k,j})$,
according to the enumeration considered in Remark~\ref{induced-srle}.

Conversely, suppose that $\tau_k(e_{k,p},e_{k,q})$ precedes
$\tau_k(e_{k,i},e_{k,j})$. If $e_{k,p}$ precedes $e_{k,i}$, then there
exists an integer $s\in\{1,\ldots ,k\}$, such that
$H_1=I_1,\ldots,H_{s-1}=I_{s-1}$ and $H_{s}$ precedes $I_{s}$. This
clearly implies that $e_{k+1,h}$ precedes $e_{k+1,l}$. On the other
hand, if $e_{k,p}=e_{k,i}$ and $e_{k,q}$ precedes $e_{k,p}$, then,
since $H_1=I_1,\ldots,H_{k}=I_{k}$, necessarily $H_k\cup H_{k+1}$
precedes $I_k\cup I_{k+1}$. But, since $H_k=I_k$, it follows that
$H_{k+1}$ precedes $I_{k+1}$ and so $e_{k+1,h}$ precedes $e_{k+1,l}$,
and $\dif_{k+1}(e_{k+1,h})$ precedes $\dif_{k+1}(e_{k+1,l})$.
\end{proof}

\begin{remark}\label{nonsrle}
If we attempt to prove Theorem~\ref{main} by means of
Theorem~\ref{algorithm}, only now using an alternative but seemingly
natural enumeration of the basis elements of the successive free
modules in the resolution, we can find that the module quotients
$M_{i}(G_{k-1})$ have an `excessive' number of generators, and as a
result, the free resolution we obtain is no longer the $\cyc$ complex.

For example, consider the $n\times n$ $\pcb$ matrix $L$ considered in
\cite[Example~1, 3]{ms}:
\begin{eqnarray*}
L=\left(\begin{array}{rrrr}
3&-1&-1&-1\\
-1&3&-1&-1\\
-1&-1&3&-1\\
-1&-1&-1&3\\
\end{array}\right).
\end{eqnarray*}
Here $\nu(L)=(1,1,1,1)$. Thus each variable $x,y,z,t$ is given degree
$1$ and we endow $A=\mbk[x,y,z,t]$ with the (weighted) reverse
lexicographic ordering (see Assumptions~\ref{grading} and
\ref{wrlo}). Now, instead of enumerating according to the $\srle$, we
enumerate $\cyc_{4,2}$ as follows (see Example~\ref{srlen=4}):
\begin{eqnarray*}
\cyc_{4,2}&=&\{(1,234),(2,134),(3,124),(123,4),(12,34),(13,24),(23,14)\}.
\end{eqnarray*}
Note that in Corollary~\ref{GBI(L)}, the enumeration considered in
$\mcb_1$ is not essential for its proof. Thus it follows that
$G_0=\{f_{0,1},f_{0,2},f_{0,3},f_{0,4},f_{0,5},f_{0,5},f_{0,6},f_{0,7}\}$
is a Gr\"obner basis of $I(\mcl)$, where
\begin{eqnarray*}
&&f_{0,1}=\dif_1(e_{1,1})=\dif_1(1,234)=\underline{x^3}-yzt,\\
&&f_{0,2}=\dif_1(e_{1,2})=\dif_1(2,134)=\underline{y^3}-xzt,\\
&&f_{0,3}=\dif_1(e_{1,3})=\dif_1(3,124)=\underline{z^3}-xyt,\\ 
&&f_{0,4}=\dif_1(e_{1,4})=\dif_1(123,4)=\underline{xyz}-t^3,\\
&&f_{0,5}=\dif_1(e_{1,5})=\dif_1(12,34)=\underline{x^2y^2}-z^2t^2,\\ 
&&f_{0,6}=\dif_1(e_{1,6})=\dif_1(13,24)=\underline{x^2z^2}-y^2t^2\mbox{ and }\\
&&f_{0,7}=\dif_1(e_{1,7})=\dif_1(23,14)=\underline{y^2z^2}-x^2t^2.
\end{eqnarray*}
It is a simple check to see that ($M_1(G_0)=0$):
\begin{eqnarray*}
&&M_2(G_0)=(x^3)\mbox{ , }
M_3(G_0)=(x^3,y^3)\mbox{ , }
M_4(G_0)=(x^2,y^2,z^2)\mbox{ , }\\
&&M_5(G_0)=M_6(G_0)=M_7(G_0)=(x,y,z)=\mfm.
\end{eqnarray*}
That gives $15$ minimal generators $x^\alpha$ of the ideal quotients
$M_i(G_0)$, for $i=2,\ldots ,7$, in Theorem~\ref{algorithm}.  Thus the
Gr\"obner basis $G_1$ of the syzygies of $G_0$ would contain $15$
elements and the corresponding resolution would begin with
\begin{eqnarray*}
0\leftarrow \mbk[x]/I(\mcl)\leftarrow
\mbk[x]\leftarrow\mbk[x]^{7}\leftarrow \mbk[x]^{15}\leftarrow\cdots ,
\end{eqnarray*}
and so it would not coincide with the complex $\mcc_{L}$.
\end{remark}

